\chapter{Methodology and Implementation}
\section{Overall pipeline}
The work runs in four stages: prepare the data, search for the best settings, retrain the chosen recipes, and test them once on held-out images. Each stage is a numbered Python script, and a few small shared modules hold the data, model, training, and optimiser code. The offline demo uses the saved results and trained models.

\ArtifactFigureShort{fig_pipeline}{System pipeline}{The project pipeline, from data preparation to search, final evaluation, and the offline demo.}{fig:pipeline}

\section{Data preparation and splitting}
We downloaded the colour tomato images from the public PlantVillage repository. The ten class folders, sorted alphabetically, define the class labels. Every image was checked to open correctly and converted to standard RGB colour. Exact duplicate images were removed. Images that look almost identical (for example, two photos of the same leaf) were grouped together and always placed in the same split, so the model is never tested on a near-copy of a training image.

The raw data had \DataOriginal\ images. We removed \DataDuplicates\ duplicates and found \DataCorrupt\ broken files, leaving \TotalN\ images. These were split about 70/15/15 into training, validation, and test sets, keeping the class proportions the same in each split. The split lists are saved in the repository so every run uses exactly the same images.

\ArtifactTable{tab_data_summary}{Number of images per class and split after cleaning.}{tab:data}

\section{Image preparation}
Images are resized once and cached to make training faster. During training, each image is randomly cropped and flipped; the searched recipes can also add RandAugment, which applies random colour and shape changes of a chosen strength. The baseline uses only cropping and flipping. Validation and test images are simply resized and centre-cropped, with no randomness.

\section{Model and baseline}
We use EfficientNet-B0 with ImageNet weights and the classifier head from Chapter~3. The network body has nine blocks. The \emph{fine-tuning depth} $k$ says how many of the last blocks are retrained: $k=0$ trains only the head, and $k=9$ retrains the whole network. The baseline trains \BaselineParams\ parameters.

\ArtifactFigureShort{fig_architecture}{EfficientNet-B0 architecture and evolved decisions}{EfficientNet-B0 with our classifier head, showing which parts the DE chromosome controls.}{fig:architecture}

\ArtifactTable{tab_model_summary}{Blocks, output shapes, and parameter counts of the baseline model.}{tab:model}

The baseline is a common ``default'' transfer-learning recipe. It was fixed before any search was run and was never changed afterwards.
\begin{table}[htbp]\centering\small
\caption{Baseline settings, fixed before the search.}
\begin{tabular}{ll}\toprule
Setting & Value\\\midrule
Backbone & ImageNet-pretrained EfficientNet-B0\\
Unfrozen feature blocks & 0\\
Classifier hidden units & 256\\
Dropout & 0.2\\
Optimizer / learning rate & Adam / $10^{-3}$\\
Weight decay / label smoothing & 0 / 0\\
Extra RandAugment & Disabled\\
Batch / final resolution & 32 / $224\times224$\\
Epoch ceiling / patience & 15 / 4\\\bottomrule
\end{tabular}
\end{table}

\section{Chromosome and fitness}
Each DE member is a \emph{chromosome} of seven numbers between 0 and 1. A decoding rule turns each number into a real setting:
\begin{table}[htbp]\centering\small
\caption{The seven genes and how each is decoded into a training setting. Whole-number settings are rounded.}
\begin{tabularx}{\textwidth}{@{}c l X@{}}\toprule
Gene & Decision & Decoded range\\\midrule
$x_1$ & Learning rate & $10^{-4.5+2.5x_1}$\\
$x_2$ & Weight decay & $10^{-6+4x_2}$\\
$x_3$ & Dropout & $0.6x_3$\\
$x_4$ & Hidden units & Index into $\{0,128,256,512\}$\\
$x_5$ & Unfrozen blocks & Integer in $\{0,\ldots,9\}$\\
$x_6$ & Label smoothing & $0.2x_6$\\
$x_7$ & RandAugment magnitude & Integer in $\{0,\ldots,15\}$\\\bottomrule
\end{tabularx}
\end{table}

The \emph{fitness} of a chromosome is the validation macro-F1 of a network trained with its settings. To keep the search affordable, each trial uses a quick training run: smaller images, part of the training set, and a few epochs. All trials use the same data and random seed, so they are compared fairly. If two chromosomes decode to the same settings, the saved score is reused. A trial that crashes receives a fitness of zero and the search continues.

The DE population starts from evenly spread random points; the baseline is not inserted into it. After the search, the best few recipes are retrained at full image size, and the one with the highest validation macro-F1 is chosen. Random search follows exactly the same rules and number of trials.

\section{Saving progress and the demo}
The search saves its progress after every trial, so it can resume after an interruption without losing work. Each trained model is saved with its settings, seed, and best epoch.

The demo loads the baseline and DE models and makes real predictions on uploaded or example leaves, with Grad-CAM heat maps. It checks uploaded files and rejects unsupported ones. After installation, it runs fully offline.
